% !TeX program = xelatex
% Báo cáo [24 - 10] — bốn mục: kết quả tuần / bộ test / kết quả test / kế hoạch.
% Số liệu: workspace/custom_suite_scores.json (chấm 24-10, mở rộng 25-10),
%         workspace/custom_suite_summary.json (mean ± SD, bootstrap 2.000 lần).
% Biên dịch: xelatex -interaction=nonstopmode "[24 - 10] - Báo cáo.tex"
\documentclass[11pt,a4paper]{article}

\usepackage{fontspec}
\setmainfont{Times New Roman}
\usepackage[a4paper,top=2.0cm,bottom=2.0cm,left=2.2cm,right=1.9cm]{geometry}
\usepackage{booktabs}
\usepackage{longtable}
\usepackage{graphicx}
\usepackage{enumitem}
\usepackage[vietnamese]{babel}
\usepackage{xcolor}

\setlist[itemize]{leftmargin=*,itemsep=1pt,topsep=2pt,parsep=0pt}
\setlength{\parindent}{0pt}
\setlength{\emergencystretch}{3em}

\title{\vspace{-1.2cm}\textbf{Báo cáo kết quả tuần 18--24/10}\\[4pt]
       {\normalsize Đề tài: Phát hiện trùng lặp mã nguồn C/C++ và Python}}
\date{\today}

\usepackage{needspace}
\usepackage{listings}
\lstset{
  basicstyle=\ttfamily\tiny,
  numbers=left,
  numberstyle=\tiny\color{gray},
  numbersep=4pt,
  frame=single,
  framesep=2pt,
  rulecolor=\color{gray},
  breaklines=true,
  columns=fullflexible,
  keepspaces=true,
  showstringspaces=false,
  tabsize=4,
  aboveskip=3pt,
  belowskip=1pt,
  xleftmargin=12pt,
}

\begin{document}
\maketitle
\vspace{-0.9cm}

%======================================================================
\section{Kết quả tuần này}
\begin{itemize}
\item Mở rộng bộ kiểm thử từ \textbf{41} lên \textbf{89 cặp}. Giữ nguyên 25 cặp
  trùng lặp và bổ sung \textbf{48 cặp không trùng lặp}, đưa nhóm này từ 16 lên
  \textbf{64 cặp} (đúng 4 lần bản cũ) cho sát tỉ lệ thực tế. Các cặp mới gồm
  24 cặp ``cùng bài, khác lõi'', 6 cặp base--base khác bài còn thiếu và 18 cặp
  khác bài nhưng cùng kiểu biến đổi (\texttt{while}/\texttt{dowhile}/
  \texttt{scanf}/\texttt{combo}). Toàn bộ sinh tất định, có kiểm tra trùng cặp.
\item Chấm \textbf{89 cặp} bằng 5 phương pháp: CodeBERT (M1), GraphCodeBERT
  (M2), ASTNN (M3), CPG-GNN (M4), mô hình kết hợp (M5). Ba phương pháp đọc văn
  bản (M1, M2, M5) nhận tối đa 512 đơn vị từ (token). MOSS (B1) chỉ có điểm
  trên \textbf{41 cặp gốc} (thiếu \texttt{MOSS\_USER\_ID}), nên mọi so sánh
  với B1 được thực hiện trên \textbf{41 cặp chung}.
\item \textbf{Kết quả chính:} cả 5 mô hình bắt đúng \textbf{25/25} cặp trùng lặp
  (recall 1,000), nhưng vẫn báo ``trùng lặp'' cho rất nhiều cặp không trùng:
  M1 23/64, M2 22/64, M5 26/64, M3 50/64, M4 62/64. F1 lần lượt: M2
  \textbf{0,694}, M1 0,685, M5 0,658, M3 0,500, M4 0,446. Khoảng cách điểm
  clone/non-clone của M3 (0,197) và M4 (0,093) quá nhỏ, nên hai mô hình này
  gần như không phân biệt được hai lớp.
\item Trên \textbf{41 cặp chung} (đủ 6 phương pháp): F1 của M2 0,909, M1 0,893,
  M5 0,877, M3 0,794, M4 0,758, B1 0,400. Như vậy thứ hạng không đổi, nhưng
  điểm tuyệt đối giảm rõ khi bộ test có nhiều cặp không trùng lặp hơn: các mô
  hình đang bám vào khung chung của hai chương trình thay vì so sánh lõi.
\end{itemize}

%======================================================================
\section{Bộ test}
\begin{center}
\small
\setlength{\tabcolsep}{6pt}
\begin{tabular}{llrrr}
\toprule
\textbf{Nhóm} & \textbf{Nội dung} & \textbf{Tệp} & \textbf{Dòng} & \textbf{Cặp} \\
\midrule
P1--P4 & 4 bài cơ bản: tổng $1..n$, ước chung lớn nhất, max mảng, số nguyên tố & 24 & 10--18 & 36 \\
P5    & Danh sách liên kết + hướng đối tượng: \texttt{struct} + hàm $\leftrightarrow$ \texttt{class} & 6 & 59--70 & 10 \\
P6    & Danh sách liên kết có thứ tự, 5 hàm, có bản \texttt{class} & 6 & 83--93 & 9 \\
Khác bài & hai bài khác đề (kiểm tra âm tính), gồm 6 cặp base--base và 18 cặp cùng biến đổi & -- & -- & 34 \\
\midrule
\textbf{Tổng} & 25 cặp trùng lặp + 64 cặp không trùng lặp & 36 & -- & \textbf{89} \\
\bottomrule
\end{tabular}
\end{center}

%======================================================================
\section{Kết quả test và đánh giá}

\subsection{Chi tiết 41 cặp gốc (đủ 6 phương pháp)}
Bảng dưới giữ nguyên 41 cặp của bộ cũ để so sánh trực tiếp với báo cáo trước.
Các ô đỏ là kết luận sai tại ngưỡng đóng băng của từng phương pháp.

{\small
\setlength{\tabcolsep}{5pt}
\begin{longtable}{llcccccc}
\toprule
\textbf{Cặp} & \textbf{Nhãn} & \textbf{B1} & \textbf{M1} & \textbf{M2} & \textbf{M3} & \textbf{M4} & \textbf{M5} \\
\midrule
\endfirsthead
\toprule
\textbf{Cặp} & \textbf{Nhãn} & \textbf{B1} & \textbf{M1} & \textbf{M2} & \textbf{M3} & \textbf{M4} & \textbf{M5} \\
\midrule
\endhead
\bottomrule
\endlastfoot
\multicolumn{8}{l}{\textbf{Nhóm P1: Tính tổng $1..n$ (5 cặp)}} \\
\midrule
P1\_while              & trùng lặp  & \textcolor{red}{0,000} & 1,000 & 1,000 & 0,945 & 0,973 & 1,000 \\
P1\_dowhile            & trùng lặp  & \textcolor{red}{0,000} & 1,000 & 1,000 & 0,760 & 0,973 & 1,000 \\
P1\_scanf              & trùng lặp  & \textcolor{red}{0,000} & 1,000 & 1,000 & 0,931 & 0,973 & 1,000 \\
P1\_combo              & trùng lặp  & \textcolor{red}{0,000} & 1,000 & 1,000 & 0,812 & 0,915 & 1,000 \\
P1\_nonclone\_fact     & khác       & 0,000 & \textcolor{red}{0,999} & \textcolor{red}{0,991} & \textcolor{red}{0,960} & \textcolor{red}{0,971} & \textcolor{red}{0,723} \\
\midrule
\multicolumn{8}{l}{\textbf{Nhóm P2: Ước chung lớn nhất (5 cặp)}} \\
\midrule
P2\_for                & trùng lặp  & \textcolor{red}{0,000} & 1,000 & 1,000 & 0,955 & 0,972 & 1,000 \\
P2\_dowhile            & trùng lặp  & \textcolor{red}{0,000} & 1,000 & 1,000 & 0,909 & 0,973 & 1,000 \\
P2\_scanf              & trùng lặp  & \textcolor{red}{0,000} & 1,000 & 1,000 & 0,834 & 0,968 & 1,000 \\
P2\_combo              & trùng lặp  & \textcolor{red}{0,000} & 1,000 & 1,000 & 0,796 & 0,973 & 1,000 \\
P2\_nonclone\_digitsum & khác       & 0,000 & 0,000 & 0,001 & \textcolor{red}{0,799} & \textcolor{red}{0,973} & 0,000 \\
\midrule
\multicolumn{8}{l}{\textbf{Nhóm P3: Tìm giá trị lớn nhất của mảng (5 cặp)}} \\
\midrule
P3\_while              & trùng lặp  & \textcolor{red}{0,000} & 1,000 & 1,000 & 0,952 & 0,973 & 1,000 \\
P3\_dowhile            & trùng lặp  & \textcolor{red}{0,000} & 1,000 & 1,000 & 0,479 & 0,973 & 1,000 \\
P3\_scanf              & trùng lặp  & \textcolor{red}{0,000} & 1,000 & 1,000 & 0,920 & 0,967 & 1,000 \\
P3\_combo              & trùng lặp  & \textcolor{red}{0,000} & 1,000 & 1,000 & 0,932 & 0,841 & 1,000 \\
P3\_nonclone\_sum      & khác       & 0,000 & 0,000 & 0,005 & \textcolor{red}{0,852} & \textcolor{red}{0,973} & 0,000 \\
\midrule
\multicolumn{8}{l}{\textbf{Nhóm P4: Kiểm tra số nguyên tố (5 cặp)}} \\
\midrule
P4\_while              & trùng lặp  & \textcolor{red}{0,000} & 1,000 & 1,000 & 0,892 & 0,967 & 1,000 \\
P4\_dowhile            & trùng lặp  & \textcolor{red}{0,000} & 1,000 & 1,000 & 0,863 & 0,967 & 1,000 \\
P4\_scanf              & trùng lặp  & \textcolor{red}{0,000} & 1,000 & 1,000 & 0,895 & 0,967 & 1,000 \\
P4\_combo              & trùng lặp  & \textcolor{red}{0,000} & 1,000 & 1,000 & 0,931 & 0,967 & 0,994 \\
P4\_nonclone\_square   & khác       & 0,000 & \textcolor{red}{1,000} & \textcolor{red}{1,000} & \textcolor{red}{0,864} & \textcolor{red}{0,967} & \textcolor{red}{1,000} \\
\midrule
\multicolumn{8}{l}{\textbf{Khác bài P1--P4: Kiểm tra âm tính (6 cặp)}} \\
\midrule
X\_P1P2                & khác       & 0,000 & 0,000 & 0,000 & \textcolor{red}{0,868} & \textcolor{red}{0,960} & 0,000 \\
X\_P1P3                & khác       & 0,000 & 0,000 & 0,174 & \textcolor{red}{0,931} & \textcolor{red}{0,973} & 0,000 \\
X\_P1P4                & khác       & 0,000 & 0,000 & 0,000 & \textcolor{red}{0,618} & \textcolor{red}{0,914} & 0,000 \\
X\_P2P3                & khác       & 0,000 & 0,000 & 0,000 & \textcolor{red}{0,913} & \textcolor{red}{0,873} & 0,000 \\
X\_P2P4                & khác       & 0,000 & 0,000 & 0,023 & \textcolor{red}{0,777} & \textcolor{red}{0,660} & \textcolor{red}{0,027} \\
X\_P3P4                & khác       & 0,000 & 0,000 & 0,000 & \textcolor{red}{0,784} & \textcolor{red}{0,660} & 0,000 \\
\midrule
\multicolumn{8}{l}{\textbf{Nhóm P5: Danh sách liên kết cơ bản + struct/class (8 cặp)}} \\
\midrule
P5\_struct\_class      & trùng lặp  & \textcolor{red}{0,000} & 1,000 & 1,000 & 0,915 & 0,973 & 1,000 \\
P5\_while              & trùng lặp  & 0,540 & 1,000 & 1,000 & 0,993 & 0,973 & 1,000 \\
P5\_scanf              & trùng lặp  & 0,710 & 1,000 & 1,000 & 0,995 & 0,973 & 1,000 \\
P5\_combo              & trùng lặp  & 0,150 & 1,000 & 1,000 & 0,994 & 0,973 & 1,000 \\
P5\_class\_combo       & trùng lặp  & \textcolor{red}{0,000} & 1,000 & 1,000 & 0,907 & 0,973 & 1,000 \\
P5\_nonclone\_palindrome & khác     & 0,000 & 0,002 & 0,000 & 0,270 & \textcolor{red}{0,743} & 0,002 \\
X\_P1P5                & khác       & 0,000 & \textcolor{red}{1,000} & \textcolor{red}{1,000} & 0,179 & \textcolor{red}{0,742} & \textcolor{red}{1,000} \\
X\_P2P5                & khác       & 0,000 & 0,000 & 0,000 & 0,375 & \textcolor{red}{0,710} & 0,000 \\
\midrule
\multicolumn{8}{l}{\textbf{Nhóm P6: Danh sách liên kết có thứ tự, $\ge 80$ dòng (7 cặp)}} \\
\midrule
P6\_struct\_class      & trùng lặp  & 0,400 & 1,000 & 1,000 & 0,970 & 0,967 & 1,000 \\
P6\_while              & trùng lặp  & 0,440 & 1,000 & 1,000 & 0,992 & 0,969 & 1,000 \\
P6\_scanf              & trùng lặp  & 0,770 & 1,000 & 1,000 & 0,996 & 0,967 & 1,000 \\
P6\_combo              & trùng lặp  & 0,200 & 1,000 & 1,000 & 0,995 & 0,967 & 1,000 \\
P6\_nonclone\_filter   & khác       & \textcolor{red}{0,400} & \textcolor{red}{1,000} & \textcolor{red}{0,998} & \textcolor{red}{0,994} & \textcolor{red}{0,973} & \textcolor{red}{0,866} \\
X\_P5P6                & khác       & \textcolor{red}{0,510} & \textcolor{red}{1,000} & \textcolor{red}{1,000} & \textcolor{red}{0,992} & \textcolor{red}{0,973} & \textcolor{red}{0,997} \\
X\_P5P6\_class         & khác       & \textcolor{red}{0,360} & \textcolor{red}{0,997} & 0,554 & \textcolor{red}{0,968} & \textcolor{red}{0,973} & \textcolor{red}{0,020} \\
\midrule
\textbf{Ngưỡng}        &            & 0,010 & 0,210 & 0,660 & 0,420 & 0,470 & 0,010 \\
\textbf{Số lỗi P1}     & 5 cặp      & 4 & 1 & 1 & 1 & 1 & 1 \\
\textbf{Số lỗi P2}     & 5 cặp      & 4 & 0 & 0 & 1 & 1 & 0 \\
\textbf{Số lỗi P3}     & 5 cặp      & 4 & 0 & 0 & 1 & 1 & 0 \\
\textbf{Số lỗi P4}     & 5 cặp      & 4 & 1 & 1 & 1 & 1 & 1 \\
\textbf{Số lỗi khác bài} & 6 cặp    & 0 & 0 & 0 & 6 & 6 & 1 \\
\textbf{Số lỗi P5}     & 8 cặp      & 2 & 1 & 1 & 0 & 3 & 1 \\
\textbf{Số lỗi P6}     & 7 cặp      & 3 & 3 & 2 & 3 & 3 & 3 \\
\midrule
\textbf{Tổng số lỗi}   & \textbf{41 cặp} & \textbf{21} & \textbf{6} & \textbf{5} & \textbf{13} & \textbf{16} & \textbf{7} \\
\end{longtable}}

\subsection{Tổng hợp mean $\pm$ SD trên toàn bộ 89 cặp}
Điểm trung bình $\pm$ SD của điểm thô đo độ phân tán giữa các cặp; F1 $\pm$ SD
là độ bất định ước lượng bằng bootstrap 2.000 lần trên các cặp (ngưỡng đóng
băng: B1 0,010; M1 0,210; M2 0,660; M3 0,420; M4 0,470; M5 0,010).

\begin{table}[htbp]\centering\small
\caption{Toàn bộ 89 cặp có điểm (M1--M5); $\bar{s}_{cl}/\bar{s}_{nc}$ là điểm
trung bình clone/non-clone $\pm$ SD}\label{tab:suite-all}
\setlength{\tabcolsep}{3pt}
\begin{tabular}{lrrrrrrrrr}
\toprule
Mô hình & $n$ & Acc & P & R & F1 & F1 $\pm$ SD & $\bar{s}_{cl}$ & $\bar{s}_{nc}$ & Khoảng cách \\
\midrule
M1 & 89 & 0,742 & 0,521 & 1,000 & 0,685 & 0,680 $\pm$ 0,064 & 1,000 $\pm$ 0,000 & 0,350 $\pm$ 0,472 & 0,650 \\
M2 & 89 & 0,753 & 0,532 & 1,000 & 0,694 & 0,689 $\pm$ 0,064 & 1,000 $\pm$ 0,000 & 0,362 $\pm$ 0,457 & 0,638 \\
M3 & 89 & 0,438 & 0,333 & 1,000 & 0,500 & 0,496 $\pm$ 0,061 & 0,903 $\pm$ 0,110 & 0,706 $\pm$ 0,263 & 0,197 \\
M4 & 89 & 0,303 & 0,287 & 1,000 & 0,446 & 0,442 $\pm$ 0,059 & 0,963 $\pm$ 0,028 & 0,870 $\pm$ 0,149 & 0,093 \\
M5 & 89 & 0,708 & 0,490 & 1,000 & 0,658 & 0,653 $\pm$ 0,064 & 1,000 $\pm$ 0,001 & 0,246 $\pm$ 0,406 & 0,753 \\
\bottomrule
\end{tabular}
\end{table}

\begin{table}[htbp]\centering\small
\caption{Điểm trung bình $\pm$ SD theo nhóm cặp (nhóm clone dùng điểm clone,
nhóm non-clone dùng điểm non-clone)}\label{tab:suite-cat-score}
\setlength{\tabcolsep}{4pt}
\begin{tabular}{lrrrr}
\toprule
Mô hình & Clone dễ & Clone khó & Không clone khó & Không clone dễ \\
\midrule
Số cặp & 16 & 9 & 30 & 34 \\
\midrule
M1 & 1,000 $\pm$ 0,000 & 1,000 $\pm$ 0,000 & 0,477 $\pm$ 0,499 & 0,238 $\pm$ 0,423 \\
M2 & 1,000 $\pm$ 0,000 & 1,000 $\pm$ 0,000 & 0,488 $\pm$ 0,491 & 0,251 $\pm$ 0,400 \\
M3 & 0,895 $\pm$ 0,128 & 0,917 $\pm$ 0,072 & 0,728 $\pm$ 0,231 & 0,686 $\pm$ 0,290 \\
M4 & 0,971 $\pm$ 0,003 & 0,950 $\pm$ 0,045 & 0,927 $\pm$ 0,094 & 0,820 $\pm$ 0,170 \\
M5 & 1,000 $\pm$ 0,000 & 0,999 $\pm$ 0,002 & 0,381 $\pm$ 0,452 & 0,127 $\pm$ 0,323 \\
\bottomrule
\end{tabular}
\end{table}

\begin{table}[htbp]\centering\small
\caption{Số lỗi (FP + FN) theo nhóm}\label{tab:suite-cat-err}
\setlength{\tabcolsep}{4pt}
\begin{tabular}{lrrrr}
\toprule
Mô hình & Clone dễ & Clone khó & Không clone khó & Không clone dễ \\
\midrule
Số cặp & 16 & 9 & 30 & 34 \\
\midrule
M1 & 0 & 0 & 15 & 8 \\
M2 & 0 & 0 & 14 & 8 \\
M3 & 0 & 0 & 26 & 24 \\
M4 & 0 & 0 & 30 & 32 \\
M5 & 0 & 0 & 16 & 10 \\
\bottomrule
\end{tabular}
\end{table}

Hai bảng theo nhóm cho thấy lỗi hoàn toàn nằm ở phía non-clone: mọi mô hình đều
đúng 25/25 cặp clone, kể cả 9 cặp ``clone khó''. Nhóm ``không clone khó'' (cùng
bài, khác lõi) khó hơn nhóm ``không clone dễ'' (khác bài) với M1, M2, M5, đúng
như kỳ vọng; riêng M3 và M4 lại cho điểm non-clone rất cao (0,706 và 0,870) nên
gần như mất khả năng phân biệt.

\subsection{So sánh công bằng trên 41 cặp chung}
\begin{table}[htbp]\centering\small
\caption{41 cặp có điểm của cả 6 phương pháp; $\bar{s}_{cl}/\bar{s}_{nc}$ là
điểm trung bình clone/non-clone $\pm$ SD}\label{tab:suite-common}
\setlength{\tabcolsep}{3pt}
\begin{tabular}{lrrrrrrrrr}
\toprule
Mô hình & $n$ & Acc & P & R & F1 & F1 $\pm$ SD & $\bar{s}_{cl}$ & $\bar{s}_{nc}$ & Khoảng cách \\
\midrule
B1 & 41 & 0,488 & 0,700 & 0,280 & 0,400 & 0,393 $\pm$ 0,106 & 0,128 $\pm$ 0,241 & 0,079 $\pm$ 0,173 & 0,049 \\
M1 & 41 & 0,854 & 0,806 & 1,000 & 0,893 & 0,890 $\pm$ 0,044 & 1,000 $\pm$ 0,000 & 0,375 $\pm$ 0,500 & 0,625 \\
M2 & 41 & 0,878 & 0,833 & 1,000 & 0,909 & 0,906 $\pm$ 0,042 & 1,000 $\pm$ 0,000 & 0,359 $\pm$ 0,466 & 0,641 \\
M3 & 41 & 0,683 & 0,658 & 1,000 & 0,794 & 0,790 $\pm$ 0,055 & 0,903 $\pm$ 0,110 & 0,759 $\pm$ 0,261 & 0,144 \\
M4 & 41 & 0,610 & 0,610 & 1,000 & 0,758 & 0,754 $\pm$ 0,058 & 0,963 $\pm$ 0,028 & 0,878 $\pm$ 0,126 & 0,086 \\
M5 & 41 & 0,829 & 0,781 & 1,000 & 0,877 & 0,874 $\pm$ 0,047 & 1,000 $\pm$ 0,001 & 0,290 $\pm$ 0,442 & 0,710 \\
\bottomrule
\end{tabular}
\end{table}

B1 thấp hơn hẳn vì MOSS chỉ nhận ra 7/25 cặp clone (recall 0,280) trên 41 cặp gốc;
điểm tuyệt đối của B1 không so sánh được với 89 cặp. Trên cùng 41 cặp, M2 vẫn
dẫn đầu (F1 0,909), tiếp theo là M1 (0,893) và M5 (0,877); khoảng cách giữa ba
mô hình này chỉ 0,032, trong khi SD bootstrap khoảng 0,04--0,05, nên chưa thể
kết luận chúng khác nhau có ý nghĩa thống kê.

%======================================================================
\section{Kế hoạch tuần tới}

%======================================================================
\section{Phụ lục: mã nguồn minh họa (P6)}

\subsection*{\texttt{p6\_base\_list.cpp} \hfill {\normalsize\mdseries 84 dòng --- bản gốc: \texttt{struct Node} + 5 hàm tự do}}
\begin{lstlisting}
#include <bits/stdc++.h>
using namespace std;

struct Node {
    int val;
    Node *next;
};

Node* insertSorted(Node *head, int v) {
    Node *node = new Node;
    node->val = v;
    node->next = NULL;
    if (head == NULL || v <= head->val) {
        node->next = head;
        return node;
    }
    Node *p = head;
    while (p->next != NULL && p->next->val <= v) {
        p = p->next;
    }
    node->next = p->next;
    p->next = node;
    return head;
}

int removeDuplicates(Node *head) {
    int removed = 0;
    Node *p = head;
    while (p != NULL && p->next != NULL) {
        if (p->val == p->next->val) {
            Node *dup = p->next;
            p->next = dup->next;
            delete dup;
            removed = removed + 1;
        } else {
            p = p->next;
        }
    }
    return removed;
}

long long sumList(Node *head) {
    long long total = 0;
    Node *p = head;
    while (p != NULL) {
        total = total + p->val;
        p = p->next;
    }
    return total;
}

void printList(Node *head) {
    Node *p = head;
    while (p != NULL) {
        cout << p->val << " ";
        p = p->next;
    }
    cout << endl;
}

void freeList(Node *head) {
    while (head != NULL) {
        Node *nxt = head->next;
        delete head;
        head = nxt;
    }
}

int main() {
    int n;
    cin >> n;
    Node *head = NULL;
    for (int i = 0; i < n; i++) {
        int x;
        cin >> x;
        head = insertSorted(head, x);
    }
    printList(head);
    cout << "removed " << removeDuplicates(head) << endl;
    printList(head);
    cout << "sum " << sumList(head) << endl;
    freeList(head);
    return 0;
}
\end{lstlisting}

\needspace{8\baselineskip}
\subsection*{\texttt{p6\_v\_while.cpp} \hfill {\normalsize\mdseries 93 dòng --- trùng lặp: \texttt{while} $\to$ \texttt{for}/\texttt{do-while}, đổi tên biến}}
\begin{lstlisting}
#include <bits/stdc++.h>
using namespace std;

struct Node {
    int val;
    Node *next;
};

Node* insertSorted(Node *head, int v) {
    Node *fresh = new Node;
    fresh->val = v;
    fresh->next = NULL;
    if (head == NULL || v <= head->val) {
        fresh->next = head;
        return fresh;
    }
    Node *cur = head;
    for (; cur->next != NULL; cur = cur->next) {
        if (cur->next->val > v) {
            break;
        }
    }
    fresh->next = cur->next;
    cur->next = fresh;
    return head;
}

int removeDuplicates(Node *head) {
    int cnt = 0;
    Node *cur = head;
    while (cur != NULL) {
        if (cur->next == NULL) {
            break;
        }
        if (cur->val != cur->next->val) {
            cur = cur->next;
            continue;
        }
        Node *dup = cur->next;
        cur->next = dup->next;
        delete dup;
        cnt = cnt + 1;
    }
    return cnt;
}

long long sumList(Node *head) {
    long long acc = 0;
    Node *cur = head;
    do {
        if (cur == NULL) {
            break;
        }
        acc = acc + cur->val;
        cur = cur->next;
    } while (cur != NULL);
    return acc;
}

void printList(Node *head) {
    for (Node *cur = head; cur != NULL; cur = cur->next) {
        cout << cur->val << " ";
    }
    cout << endl;
}

void freeList(Node *head) {
    Node *cur = head;
    while (cur != NULL) {
        Node *nxt = cur->next;
        delete cur;
        cur = nxt;
    }
}

int main() {
    int n;
    cin >> n;
    Node *head = NULL;
    int i = 0;
    while (i < n) {
        int x;
        cin >> x;
        head = insertSorted(head, x);
        i = i + 1;
    }
    printList(head);
    cout << "removed " << removeDuplicates(head) << endl;
    printList(head);
    cout << "sum " << sumList(head) << endl;
    freeList(head);
    return 0;
}
\end{lstlisting}

\needspace{8\baselineskip}
\subsection*{\texttt{p6\_v\_scanf.cpp} \hfill {\normalsize\mdseries 83 dòng --- trùng lặp: \texttt{cin/cout} $\to$ \texttt{scanf/printf}}}
\begin{lstlisting}
#include <cstdio>

struct Node {
    int val;
    Node *next;
};

Node* insertSorted(Node *head, int v) {
    Node *node = new Node;
    node->val = v;
    node->next = NULL;
    if (head == NULL || v <= head->val) {
        node->next = head;
        return node;
    }
    Node *p = head;
    while (p->next != NULL && p->next->val <= v) {
        p = p->next;
    }
    node->next = p->next;
    p->next = node;
    return head;
}

int removeDuplicates(Node *head) {
    int removed = 0;
    Node *p = head;
    while (p != NULL && p->next != NULL) {
        if (p->val == p->next->val) {
            Node *dup = p->next;
            p->next = dup->next;
            delete dup;
            removed = removed + 1;
        } else {
            p = p->next;
        }
    }
    return removed;
}

long long sumList(Node *head) {
    long long total = 0;
    Node *p = head;
    while (p != NULL) {
        total = total + p->val;
        p = p->next;
    }
    return total;
}

void printList(Node *head) {
    Node *p = head;
    while (p != NULL) {
        printf("%d ", p->val);
        p = p->next;
    }
    printf("\n");
}

void freeList(Node *head) {
    while (head != NULL) {
        Node *nxt = head->next;
        delete head;
        head = nxt;
    }
}

int main() {
    int n;
    scanf("%d", &n);
    Node *head = NULL;
    for (int i = 0; i < n; i++) {
        int x;
        scanf("%d", &x);
        head = insertSorted(head, x);
    }
    printList(head);
    printf("removed %d\n", removeDuplicates(head));
    printList(head);
    printf("sum %lld\n", sumList(head));
    freeList(head);
    return 0;
}
\end{lstlisting}

\needspace{8\baselineskip}
\subsection*{\texttt{p6\_v\_combo.cpp} \hfill {\normalsize\mdseries 87 dòng --- trùng lặp khó: đảo điều kiện + biến thừa}}
\begin{lstlisting}
#include <cstdio>

struct Node {
    int val;
    Node *next;
};

Node* addNode(Node *head, int v) {
    Node *fresh = new Node;
    fresh->val = v;
    fresh->next = NULL;
    if (!(head != NULL && v > head->val)) {
        fresh->next = head;
        return fresh;
    }
    Node *cur = head;
    while (!(cur->next == NULL || cur->next->val > v)) {
        cur = cur->next;
    }
    fresh->next = cur->next;
    cur->next = fresh;
    return head;
}

int dropRepeat(Node *head) {
    int cnt = 0;
    int unused = 7;
    Node *cur = head;
    do {
        if (cur == NULL || cur->next == NULL) {
            break;
        }
        if (cur->val != cur->next->val) {
            cur = cur->next;
            continue;
        }
        Node *dup = cur->next;
        cur->next = dup->next;
        delete dup;
        cnt = cnt + 1;
    } while (cur != NULL);
    return cnt + unused * 0;
}

long long accSum(Node *head) {
    long long acc = 0;
    Node *cur = head;
    while (!(cur == NULL)) {
        acc = acc + cur->val;
        cur = cur->next;
    }
    return acc;
}

void show(Node *head) {
    for (Node *cur = head; cur != NULL; cur = cur->next) {
        printf("%d ", cur->val);
    }
    printf("\n");
}

void release(Node *head) {
    while (!(head == NULL)) {
        Node *nxt = head->next;
        delete head;
        head = nxt;
    }
}

int main() {
    int n;
    scanf("%d", &n);
    Node *head = NULL;
    int i = 0;
    while (i < n) {
        int x;
        scanf("%d", &x);
        head = addNode(head, x);
        i = i + 1;
    }
    show(head);
    printf("removed %d\n", dropRepeat(head));
    show(head);
    printf("sum %lld\n", accSum(head));
    release(head);
    return 0;
}
\end{lstlisting}

\needspace{8\baselineskip}
\subsection*{\texttt{p6\_v\_class.cpp} \hfill {\normalsize\mdseries 86 dòng --- trùng lặp: \texttt{class IntList} + hàm hủy, cùng hành vi}}
\begin{lstlisting}
#include <bits/stdc++.h>
using namespace std;

struct Node {
    int val;
    Node *next;
};

class IntList {
private:
    Node *head;
public:
    IntList() {
        head = NULL;
    }
    ~IntList() {
        clear();
    }
    void insert(int v) {
        Node *node = new Node;
        node->val = v;
        node->next = NULL;
        if (head == NULL || v <= head->val) {
            node->next = head;
            head = node;
            return;
        }
        Node *p = head;
        while (p->next != NULL && p->next->val <= v) {
            p = p->next;
        }
        node->next = p->next;
        p->next = node;
    }
    int dedup() {
        int removed = 0;
        Node *p = head;
        while (p != NULL && p->next != NULL) {
            if (p->val == p->next->val) {
                Node *dup = p->next;
                p->next = dup->next;
                delete dup;
                removed = removed + 1;
            } else {
                p = p->next;
            }
        }
        return removed;
    }
    long long total() const {
        long long sum = 0;
        for (Node *p = head; p != NULL; p = p->next) {
            sum = sum + p->val;
        }
        return sum;
    }
    void show() const {
        for (Node *p = head; p != NULL; p = p->next) {
            cout << p->val << " ";
        }
        cout << endl;
    }
    void clear() {
        while (head != NULL) {
            Node *nxt = head->next;
            delete head;
            head = nxt;
        }
    }
};

int main() {
    int n;
    cin >> n;
    IntList list;
    for (int i = 0; i < n; i++) {
        int x;
        cin >> x;
        list.insert(x);
    }
    list.show();
    cout << "removed " << list.dedup() << endl;
    list.show();
    cout << "sum " << list.total() << endl;
    return 0;
}
\end{lstlisting}

\needspace{8\baselineskip}
\subsection*{\texttt{p6\_n\_filter.cpp} \hfill {\normalsize\mdseries 88 dòng --- không trùng lặp: cùng khung, khác lõi ở giữa tệp}}
\begin{lstlisting}
#include <bits/stdc++.h>
using namespace std;

struct Node {
    int val;
    Node *next;
};

Node* pushBack(Node *head, int v) {
    Node *node = new Node;
    node->val = v;
    node->next = NULL;
    if (head == NULL) {
        return node;
    }
    Node *p = head;
    while (p->next != NULL) {
        p = p->next;
    }
    p->next = node;
    return head;
}

int removeGreater(Node *&head, int limit) {
    int removed = 0;
    while (head != NULL && head->val > limit) {
        Node *old = head;
        head = head->next;
        delete old;
        removed = removed + 1;
    }
    Node *p = head;
    while (p != NULL && p->next != NULL) {
        if (p->next->val > limit) {
            Node *old = p->next;
            p->next = old->next;
            delete old;
            removed = removed + 1;
        } else {
            p = p->next;
        }
    }
    return removed;
}

int countOdd(Node *head) {
    int odd = 0;
    for (Node *p = head; p != NULL; p = p->next) {
        if (p->val % 2 != 0) {
            odd = odd + 1;
        }
    }
    return odd;
}

void printList(Node *head) {
    for (Node *p = head; p != NULL; p = p->next) {
        cout << p->val << " ";
    }
    cout << endl;
}

void freeList(Node *head) {
    while (head != NULL) {
        Node *nxt = head->next;
        delete head;
        head = nxt;
    }
}

int main() {
    int n;
    cin >> n;
    Node *head = NULL;
    for (int i = 0; i < n; i++) {
        int x;
        cin >> x;
        head = pushBack(head, x);
    }
    printList(head);
    int limit;
    cin >> limit;
    cout << "removed " << removeGreater(head, limit) << endl;
    printList(head);
    cout << "odd " << countOdd(head) << endl;
    freeList(head);
    return 0;
}
\end{lstlisting}

\end{document}
